\documentclass[10pt,a4paper]{article}

% ---------- page setup (no geometry package needed) ----------
\setlength{\paperwidth}{210mm}\setlength{\paperheight}{297mm}
\setlength{\pdfpagewidth}{210mm}\setlength{\pdfpageheight}{297mm}
\setlength{\textwidth}{186mm}\setlength{\oddsidemargin}{-13.4mm}
\setlength{\textheight}{267mm}\setlength{\topmargin}{-19mm}
\setlength{\headheight}{0pt}\setlength{\headsep}{0pt}\setlength{\footskip}{0pt}
\setlength{\parindent}{0pt}
\pagestyle{empty}

\usepackage[hidelinks]{hyperref}

% ---------- helpers ----------
\renewcommand{\section}[1]{\vspace{6pt}{\large\bfseries\MakeUppercase{#1}}\\[-7pt]\rule{\textwidth}{0.6pt}\vspace{2pt}}
\newcommand{\entry}[2]{\par\textbf{#1}\hfill #2\par}
\newenvironment{points}{\vspace{-4pt}\begin{itemize}\setlength{\itemsep}{0pt}\setlength{\parskip}{0pt}\setlength{\topsep}{0pt}\setlength{\partopsep}{0pt}\setlength{\leftmargin}{12pt}}{\end{itemize}\vspace{-2pt}}
\renewcommand{\labelitemi}{$\bullet$}
\newcommand{\sep}{\enspace$|$\enspace}

\begin{document}

% ---------- header ----------
\begin{center}
  {\Huge\bfseries ADARSH MILAN}\\[3pt]
  {\large Android \,---\, Mobile \& Game Developer}\\[3pt]
  Gujarat, India \sep
  \href{mailto:adarshmilan92@gmail.com}{adarshmilan92@gmail.com} \sep
  +91 6354521829\\[1pt]
  \href{https://www.linkedin.com/in/adarsh-milan-383527304}{linkedin.com/in/adarsh-milan-383527304} \sep
  \href{https://github.com/Adarshmilan}{github.com/Adarshmilan} \sep
  \href{https://portfoliyo-zeta-indol.vercel.app/}{portfoliyo-zeta-indol.vercel.app}
\end{center}

Mobile developer and B.E.\ Computer Science student who builds Android apps and games. My games are
published on Google Play under the \textbf{2Hex} developer account, and my mobile-first web apps serve \textbf{400+
real users}. Experienced with Kotlin, Android's VpnService, Retrofit and Coroutines, and with Godot 4 for games,
plus the backend skills to build complete apps.

\section{Education}
\entry{Marwadi University, Rajkot}{Aug 2024 -- Dec 2028}
Bachelor of Engineering --- Computer Science \hfill GPA: 8.5\\

\section{Skills}
\textbf{Languages:} Kotlin, Java, GDScript, Python, JavaScript, TypeScript, SQL\\
\textbf{Android:} Android SDK, Material 3, ViewPager2, VpnService, Retrofit, OkHttp, Coroutines, Coil, Gradle (Kotlin DSL)\\
\textbf{Game Development:} Godot 4, GDScript, Shaders (toon, outline), 3D model integration, Android export\\
\textbf{Backend \& Web:} Node.js, Express, Supabase, REST APIs, React, PWA\\
\textbf{Tools:} Android Studio, Git, GitHub, Linux, Docker\\

\section{Experience}
\entry{Northstar Innovation --- Software Engineering Intern (Remote)}{Aug 2026 -- Present}
\begin{points}
  \item Selected through a technical evaluation that involved building a Shopify Admin GraphQL API $\rightarrow$ PostgreSQL product data pipeline with pagination, upserts and error handling.
  \item Work on backend system design and database management.
  \item Integrate AI tools into software development and documentation workflows.
\end{points}

\section{Projects}
\entry{GateBreak --- Mobile Game \textnormal{\sep Godot 4, GDScript, Shaders, Android}}{Google Play (in review)}
\begin{points}
  \item Developed a merge tower-defense game for Android: drag-and-merge turret grid, wave director with boss waves, and endless, data-driven difficulty scaling.
  \item Rendered 3D turret models inside a 2D game with custom toon and outline shaders; added a portrait layout that fits any screen, haptics and a JSON save system.
  \item Published under the \textbf{2Hex} Google Play developer account (GateBreak currently in review).
\end{points}
\entry{PhoneShark --- No-Root Packet Analyzer for Android \textnormal{\sep Kotlin, VpnService, Java NIO}}{}
\begin{points}
  \item Built a Wireshark-style app that captures all of a phone's traffic without root, using a local VPN and a user-space TCP/UDP forwarding engine.
  \item Wrote decoders for IPv4/IPv6, TCP, UDP, ICMP, DNS, HTTP, TLS (SNI, ALPN) and QUIC; per-app traffic attribution; \texttt{.pcap} export; unit tests.
\end{points}
\entry{Epoch --- ``On This Day'' Android App \textnormal{\sep Kotlin, Retrofit, Coroutines, Coil}}{}
\begin{points}
  \item Built a swipeable Android app (ViewPager2) that loads today's historical events from the Wikipedia REST API with Retrofit and Coroutines, with images via Coil.
\end{points}
\entry{Cyborg City --- Esports Tournament Platform (PWA) \textnormal{\sep React, Vite PWA, Supabase}}{\href{https://cyborgcity.vercel.app}{cyborgcity.vercel.app}}
\begin{points}
  \item Built a mobile-first, installable PWA tournament platform used by \textbf{200+ players}, with team management, match registration, Razorpay payments and a live leaderboard.
\end{points}
\entry{UniCart --- Campus Marketplace}{}
\begin{points}
  \item Built a campus delivery marketplace where students post requests and peers deliver them for a fee. It has \textbf{200+ real users} and is still growing.
\end{points}

\entry{Autouploader --- AI Shorts \& Reels Generator \textnormal{\sep Python, FFmpeg, yt-dlp, OpenRouter, Docker}}{\href{https://github.com/Adarshmilan/autouploader}{GitHub}}
\begin{points}
  \item Built a self-hosted web app that turns a long YouTube video into vertical Shorts/Reels: it pulls the transcript, has an LLM pick the most hyped moments, downloads only those parts and cuts 1080$\times$1920 clips with FFmpeg.
  \item Posts reviewed clips to YouTube (OAuth) and Instagram (Graph API) with a background worker that tracks each platforms daily limit and queues clips over the cap; runs in Docker.
\end{points}
\entry{Unlearning in Continual Learning Is Shallow --- First-Author Paper \textnormal{\sep Python, PyTorch, ResNet-18}}{\href{https://github.com/Adarshmilan/CUL-continual-unlearning-}{GitHub}}
\begin{points}
  \item Showed that machine unlearning in continual learning is shallow: after a forget request the output is at chance, yet linear probes on every ResNet-18 block still recover the forgotten task, and it comes back after the model learns one new task.
  \item Ran 114 experiments ($\sim$30 GPU hours) across 5 unlearning methods (BadTeacher, NegGrad+, SCRUB, SalUn, random labels), CIFAR-100 and Tiny-ImageNet, and the DER++ and ER replay learners; released the code and a reference-aligned evaluation protocol.
\end{points}

\section{Achievements \& Certifications}
\begin{points}
  \item Smart India Hackathon (SIH): built the \textit{Swachhta Sahayak} smart waste-management prototype.
  \item TCS Hackathon (\textit{CourseWise AI}) \sep Hacktoberfest 2026 (\textit{GitCompass 404}, Team 404 Not Found).
  \item \textbf{Google (Coursera):} Google AI Essentials (Oct 2024) \sep Foundations of Cybersecurity (Sep 2024) \sep Play It Safe: Manage Security Risks (Nov 2024) \sep Tools of the Trade: Linux and SQL (Jan 2025) \sep Connect and Protect: Networks and Network Security (Mar 2025).
  \item \textbf{Cisco Networking Academy:} CCNA: Introduction to Networks (Feb 2026) \sep CCNA: Switching, Routing, and Wireless Essentials (Mar 2026) \sep Linux Essentials (Cisco \& NDG, Mar 2026) \sep Operating Systems Basics (Feb 2026).
\end{points}

\section{Additional Information}
\textbf{Languages:} English, Hindi\\
\textbf{Strengths:} Product thinking, UI polish, Fast learner, Problem solving, Teamwork
\end{document}
